\section{Physics fundamentals, from scratch}
\label{sec:physics}

This section assumes you know nothing about heat physics. Take it slowly.

\subsection{Heat vs.\ temperature (they are not the same thing)}

\begin{itemize}[leftmargin=1.4em,itemsep=3pt]
  \item \textbf{Temperature} is \emph{how hot something is} --- a single number
    at a point, like \SI{300}{\kelvin}. It tells you the average jiggling energy
    of the atoms there.
  \item \textbf{Heat} is \emph{energy in transit} --- energy flowing from a
    hotter place to a colder place. You measure it as a \emph{rate of flow}.
\end{itemize}

\begin{analogy}
Temperature is like the \textbf{water level} in a tank; heat is like the
\textbf{water flowing} between tanks. Two tanks at the same level (same
temperature) have no flow between them, even if one is huge and one is tiny.
\end{analogy}

\textbf{Kelvin (K)} is just the temperature scale scientists use. It starts at
absolute zero (the coldest possible), so there are no negative numbers.
\SI{0}{\celsius} $=$ \SI{273.15}{\kelvin}. Lunar daytime soil-surface can hit
${\sim}\SI{390}{\kelvin}$ (scorching); lunar night can fall below
\SI{100}{\kelvin} ($-\SI{170}{\celsius}$). A metre down, it settles to a
near-constant ${\sim}\SI{250}{\kelvin}$.

\subsection{Conduction: how heat moves through a solid}

There are three ways heat travels: \textbf{conduction} (through direct contact,
atom bumping atom), \textbf{convection} (carried by a moving fluid like air or
water), and \textbf{radiation} (electromagnetic waves, like sunlight or the
warmth you feel from a fire).

The Moon has \textbf{no air and no water underground}, so inside the regolith
column there is essentially only \textbf{conduction} (plus a small radiative
effect across the tiny gaps between dust grains, which we fold into an
``effective'' conductivity --- see \autoref{sec:conductivity}). This is why the
model is ``1-D conduction.''

The basic law of conduction (\textbf{Fourier's law}) says, in words: \emph{heat
flows from hot to cold, and the rate of flow is proportional to how steep the
temperature change is, times how good a conductor the material is.} As a tiny
equation,
\begin{equation}
  q \;=\; -\,K\,\frac{\mathrm{d}T}{\mathrm{d}z},
  \label{eq:fourier}
\end{equation}
where:
\begin{itemize}[leftmargin=1.4em,itemsep=3pt]
  \item $\mathrm{d}T/\mathrm{d}z$ is the \textbf{temperature gradient}: how fast
    temperature changes as you go down ($z$ is depth). Steep gradient $\to$ fast
    flow.
  \item $K$ is the \textbf{thermal conductivity}: the material property --- high
    $K$ means heat flows easily (metal), low $K$ means it resists (styrofoam,
    lunar dust).
  \item The minus sign just means heat flows \emph{down} the gradient, from hot
    to cold.
\end{itemize}
In this repo Fourier's law literally appears as the conductive term in the
surface energy balance and in every interior step. See
\code{surface\_energy\_balance\_residual} in \file{lunar/solver.py}.

\subsection{The 1-D heat equation --- full derivation from scratch}
\label{sec:heateq}

This section answers three questions in order: (1)~what problem is the code
actually solving? (2)~where does the differential equation come from, step by
step, with no prior PDE course assumed? (3)~what do the top and bottom
boundaries mean? The final equation is the one in \file{lunar/solver.py} and
in the paper (Eq.~\eqref{eq:heat}).

\subsubsection{The problem in one sentence}

\begin{keyidea}
Given how sunlight hits the lunar surface hour by hour, and given how the soil
conducts heat at each depth, \textbf{find} the temperature $T$ at every depth
$z$ and every time $t$.
\end{keyidea}

The simulator predicts a \textbf{temperature profile through the soil column}
--- temperatures from the surface down to \SI{5}{m} --- and updates that
profile as the lunar day/night cycle rolls on. In code, the answer lives in
\code{PixelOutputs.T} from \code{solve\_pixel} in \file{lunar/solver.py}: a
2-D array of shape \code{(N\_z, N\_t)}.

\subsubsection{Picture a thin horizontal slice of soil}

Imagine the regolith as a \textbf{stack of thin horizontal pancakes}, each
between depth $z$ and $z+\Delta z$. Focus on one pancake in the middle:

\begin{verbatim}
     surface (z = 0)  <- sunlight, radiation to space
  -----------------------------------------
        warmer layer
  -----------------------------------------   <- top face (depth z)
        OUR SLICE  (thickness Delta z)
  -----------------------------------------   <- bottom face (z + Delta z)
        cooler layer
  -----------------------------------------
        ...
  -----------------------------------------
     bottom (z = 5 m)  <- steady heat from interior
\end{verbatim}

We track a column of soil \SI{1}{m} $\times$ \SI{1}{m} in footprint. The
slice has mass $\rho\,\Delta z$ per m$^2$ and thermal mass $\rho\,c_p\,\Delta z$
--- energy needed to warm it by \SI{1}{\kelvin}. The product $\rho\,c_p$ is
\textbf{thermal inertia}: dense, high-$c_p$ soil is hard to warm or cool (like a
heavy cast-iron pan).

\subsubsection{Heat flux at each face}

\textbf{Heat flux} $q$ [\unit{W.m^{-2}}] is the rate thermal energy crosses a
horizontal plane. Fourier's law (Eq.~\eqref{eq:fourier}) gives
\begin{equation}
  q(z) \;=\; -\,K(T,z)\,\frac{\partial T}{\partial z}.
\end{equation}
Heat flows from hot toward cold. For a slice between $z$ and $z+\Delta z$,
let $q_\mathrm{in}$ be flux across the top face and $q_\mathrm{out}$ across the
bottom.

\subsubsection{Energy conservation (the bookkeeping step)}

Energy is not created or destroyed in the slice:
\begin{equation}
  \underbrace{\rho\,c_p\,\Delta z\,\frac{\partial T}{\partial t}}_{\text{rate stored energy changes}}
  \;=\;
  q_\mathrm{in} - q_\mathrm{out}.
  \label{eq:slice-balance}
\end{equation}
The left side is how fast the slice warms; the right side is net heat flowing
in.

\subsubsection{From two faces to a derivative}

For a thin slice, $q_\mathrm{out} - q_\mathrm{in} \approx
(\partial q/\partial z)\,\Delta z$, so
$q_\mathrm{in} - q_\mathrm{out} \approx -(\partial q/\partial z)\,\Delta z$.
Substitute into Eq.~\eqref{eq:slice-balance} and cancel $\Delta z$:
\begin{equation}
  \rho\,c_p\,\frac{\partial T}{\partial t} \;=\; -\,\frac{\partial q}{\partial z}.
\end{equation}
Insert $q = -K\,\partial T/\partial z$:
\begin{equation}
  \rho(z)\,c_p(T)\,\frac{\partial T}{\partial t}
   \;=\; \frac{\partial}{\partial z}\!\left[\,K(T,z)\,\frac{\partial T}{\partial z}\,\right].
  \label{eq:heat}
\end{equation}
That is the \textbf{1-D heat equation} --- the single rule the whole simulator
enforces.

\begin{plainenglish}
\textbf{Why ``partial'' derivatives?} $T$ depends on \emph{both} depth $z$
\emph{and} time $t$. $\partial T/\partial t$ means ``how fast temperature at
this depth is changing right now''; $\partial T/\partial z$ means ``how steep
the profile is at this instant.'' You do not need a full PDE course to use
this --- treat them as separate rates of change.
\end{plainenglish}

\begin{plainenglish}
\textbf{Discrete version (what the computer does).} On a finite grid,
\begin{equation*}
  (\rho c_p)_i\,\frac{T_i^{n+1}-T_i^{n}}{\Delta t}
  = \text{(flux in from above)} - \text{(flux out below)}.
\end{equation*}
Crank--Nicolson (\autoref{sec:cn}) rearranges this so $\Delta t =
\SI{1}{h}$ is stable. Inside \code{\_step} in \file{lunar/solver.py}:
\code{cap\_i * dT\_i/dt = flux\_right(i) - flux\_left(i)}.
\end{plainenglish}

\subsubsection{Reading the final equation, term by term}

\begin{center}
\small
\begin{tabular}{@{}lll@{}}
\toprule
Piece & Plain English & Code \\
\midrule
$T(z,t)$ & Temperature at depth $z$, time $t$ & \code{PixelOutputs.T[i,k]} \\
$\partial T/\partial t$ & Warming/cooling rate & finite difference over $\Delta t$ \\
$\rho(z)$, $c_p(T)$ & Density, specific heat & \code{rho\_func}, \code{cp\_func} \\
$\rho\,c_p$ & Thermal inertia & \code{cap} in \code{\_step} \\
$K(T,z)$ & Conductivity & \code{K\_func} \\
$K\,\partial T/\partial z$ & Heat flux (Fourier) & face fluxes in \code{\_step} \\
$\partial/\partial z[\cdots]$ & Net flux in minus out & difference of face fluxes \\
\bottomrule
\end{tabular}
\end{center}

\textbf{Left side} = storage; \textbf{right side} = net conduction from
neighbours. When $K$ varies with $z$ or $T$, it stays \emph{inside} the
$\partial/\partial z$ --- the same gradient can carry different flux where
conductivity changes.

\subsubsection{The two boundaries (summary)}

The PDE needs two boundary conditions. \textbf{\autoref{sec:bc-full}} derives
them from scratch; \textbf{\autoref{sec:orbit-top-bc}} adds orbital geometry,
insolation diagrams, and the surface energy-balance figure.
both from scratch and explains their connection to the interior equation and
to \file{lunar/solver.py}. In brief: top radiative balance
(Eq.~\eqref{eq:bc-top-teach}); bottom geothermal flux
(Eq.~\eqref{eq:bc-bot-teach}).

\subsubsection{What ``solving'' means in this repository}

The simulator solves an \textbf{initial-boundary-value problem}:
\begin{enumerate}[leftmargin=1.5em,itemsep=3pt]
  \item \textbf{PDE} (interior): Eq.~\eqref{eq:heat}.
  \item \textbf{Top BC} (each timestep): Eq.~\eqref{eq:bc-top-teach}.
  \item \textbf{Bottom BC}: Eq.~\eqref{eq:bc-bot-teach}.
  \item \textbf{Initial condition}: starting profile $T(z,t{=}0)$ --- the
    equilibrium driver (\autoref{sec:anchor}) removes dependence on the guess.
\end{enumerate}

\textbf{One timestep} in \code{solve\_pixel}: evaluate $K,\rho,c_p$; assemble
Crank--Nicolson; Newton solve for $T_s$; tridiagonal solve (\code{\_thomas}) for
interior cells; advance by $\Delta t$. Repeat $\sim$709 steps per lunation.

\begin{headsup}
\textbf{Sanity checks.} (i)~Flat $T$ $\Rightarrow$ zero flux $\Rightarrow$
$\partial T/\partial t = 0$. (ii)~Steady state with bottom heat $\Rightarrow$
$\mathrm{d}\langle T\rangle/\mathrm{d}z = Q_b/K$ (geothermal gradient).
(iii)~Larger $K$ $\Rightarrow$ deeper diurnal damping (larger skin depth).
These are exactly what the code exhibits.
\end{headsup}

``1-D'' means we only track depth ($z$), ignoring sideways variation --- valid
over the small patch around a borehole.

\subsection{The three soil properties --- what they mean and why the heat equation needs them}
\label{sec:properties}

\subsubsection{Thermal capacity: why $\rho\,c_p$ sits on the left}

When you warm soil, you store \textbf{internal energy}. For a slice of thickness
$\Delta z$ and footprint \SI{1}{m^2}:
\[
  \text{mass} = \rho\,\Delta z \quad (\si{kg/m^2}), \qquad
  \text{energy to warm by \SI{1}{\kelvin}} = \rho\,c_p\,\Delta z \quad (\si{J/m^2}).
\]
So $\rho\,c_p$ is \textbf{volumetric heat capacity} (\si{J.m^{-3}.K^{-1}}). Dividing
the energy-balance rate by this capacity gives the warming rate --- exactly why
$\rho$ and $c_p$ multiply on the \textbf{left} of Eq.~\eqref{eq:heat}.

\begin{center}
\small
\begin{tabular}{@{}llll@{}}
\toprule
Symbol & Name & Units & Code \\
\midrule
$\rho(z)$ & Density & \si{kg/m^3} & \code{density\_hayne} \\
$c_p(T)$ & Specific heat & \si{J.kg^{-1}.K^{-1}} & \code{specific\_heat} \\
$\rho c_p$ & Volumetric heat capacity & \si{J.m^{-3}.K^{-1}} & \code{cap} in \code{\_step} \\
$K(T,z)$ & Conductivity & \si{W.m^{-1}.K^{-1}} & \code{conductivity\_hayne} \\
\bottomrule
\end{tabular}
\end{center}

\begin{analogy}
$\rho\,c_p$ is the soil's \textbf{thermal mass} (hard to change temperature).
$K$ is its \textbf{conductance} (how easily heat travels through).
\end{analogy}

All three are evaluated at every grid cell from \file{lunar/properties.py}.

\subsection{Thermal conductivity \texorpdfstring{$K(T,z)$}{K(T,z)} --- full physical picture}
\label{sec:conductivity}

\Kd{} --- the quantity this project retrieves --- lives inside the conductivity
formula. Two channels carry heat in vacuum regolith:

\begin{enumerate}[leftmargin=1.5em,itemsep=3pt]
  \item \textbf{Contact conduction} --- grain-to-grain; weak when fluffy.
  \item \textbf{Radiative conduction across vacuum gaps} --- grows with $T^3$.
\end{enumerate}

Hayne splits these explicitly:
\[
  K(T,z) = \underbrace{K_c(z)}_{\text{structural}} \times
           \underbrace{\bigl[1 + \chi\,(T/350)^3\bigr]}_{\text{radiative boost}}.
\]

\subsubsection{Structural part $K_c(z)$ --- compaction with depth}
\begin{equation}
  K_c(z) = K_d - (K_d - K_s)\,e^{-z/H}.
  \label{eq:kc}
\end{equation}
At $z{=}0$: $K_c{=}K_s$ (fluffy surface, ${\sim}0.74$~\mwmk). Deep down:
$K_c\to K_d$ (the unknown we fit). The scale height $H\sim\SI{6}{cm}$ sets how
fast the transition happens; at $z{=}H$ the rise is ${\sim}63\%$ complete.

\subsubsection{Radiative boost --- why $T^3$?}
Across a vacuum gap, net radiative flux $\propto T_1^4 - T_2^4 \approx 4T^3\Delta T$
for small $\Delta T$. Hayne folds this into:
\begin{equation}
  K(T,z) = K_c(z)\,\bigl[\,1 + \chi\,(T/350)^3\,\bigr].
  \label{eq:krad}
\end{equation}
At $T\approx\SI{250}{\kelvin}$ the factor is ${\sim}2$ --- radiation roughly
\textbf{doubles} effective $K$. \textbf{Not optional} (known bug if omitted).

\subsubsection{Density and specific heat}
Density compacts with the same exponential form (Hayne Eq.~5):
$\rho(z)=\rho_d-(\rho_d-\rho_s)e^{-z/H}$.
Specific heat $c_p(T)$ is a 4th-order polynomial (\code{specific\_heat}); it sets
thermal inertia but is not the retrieval's free parameter.

So \Kd{} is one knob in a fixed published shape; the retrieval freezes $K_s$, $H$,
$\chi$, and slides only \Kd{} until the model matches Apollo data.

\begin{lstlisting}[caption={Cited conductivity constants in \file{lunar/constants.py}.}]
# Surface contact conductivity [W m^-1 K^-1] -- Hayne et al. (2017) Table 2.
K_SURFACE: float = 7.4e-4
# Deep contact conductivity [W m^-1 K^-1] -- Hayne et al. (2017) Table 2.
K_DEEP: float = 3.4e-3
# Radiative conductivity coefficient chi [dimensionless] -- Hayne et al. (2017).
CHI_RADIATIVE: float = 2.7
\end{lstlisting}

\begin{headsup}
\textbf{Units.} Conductivity is in \unit{W.m^{-1}.K^{-1}} (watts per metre per
kelvin). $3.4\times10^{-3}$~\unit{W.m^{-1}.K^{-1}} $=$ \num{3.4}~\mwmk{} --- the
paper quotes milliwatts (mW) because the numbers are tiny. \SI{1}{W} $=$
\SI{1000}{mW}. So the committed \Kdstar{}~$=4.58\times10^{-3}$ in the JSON prints
as \num{4.58}~\mwmk{}.
\end{headsup}

\begin{figure}[h]
\centering
\includegraphics[width=0.78\textwidth]{fig_conductivity.png}
\caption[Conductivity rises with depth and temperature]{The Hayne $K(T,z)$ law,
Eqs.~\eqref{eq:kc}--\eqref{eq:krad}, at the Apollo~15 \Kdstar. \textbf{What to
notice:} each curve climbs steeply over the first few centimetres --- that is the
structural rise from fluffy surface ($K_s$) toward the deep value $K_d$ governed
by the scale height $H$. The three curves (cold to hot) are spread apart by the
radiative $T^3$ boost: hotter soil conducts \emph{better} even at the same depth.
\textbf{Why it matters:} the single deep number we retrieve, \Kdstar, is the
plateau these curves are climbing toward.}
\label{fig:conductivity}
\end{figure}

\subsection{The thermal skin and skin depth --- full derivation}
\label{sec:skin}

\subsubsection{Thermal diffusivity}
\[
  \alpha = \frac{K}{\rho\,c_p} \quad [\si{m^2/s}].
\]
Dividing the heat equation by $\rho c_p$ (constant properties):
$\partial T/\partial t = \alpha\,\partial^2 T/\partial z^2$ --- a diffusion equation.

\subsubsection{The decaying thermal wave}
For sinusoidal surface forcing $T(0,t)=T_{\mathrm{mean}}+A_s\cos(\omega t)$:
\[
  T(z,t) = T_{\mathrm{mean}} + A_s\,e^{-z/\delta}\cos(\omega t - z/\delta).
\]
At $z=\delta$ the amplitude is $A_s/e\approx 0.37A_s$.

\subsubsection{Deriving the skin depth}
Substituting into the diffusion equation gives
\begin{equation}
  \delta = \sqrt{\frac{2\alpha}{\omega}}.
  \label{eq:skin}
\end{equation}
Larger $\alpha$ $\Rightarrow$ deeper skin; larger $\omega$ $\Rightarrow$ shallower
skin. For the Moon, $\delta\sim\SI{6}{cm}$.

Apollo sensors at \SI{0.8}{m}--\SI{2.4}{m} sit \textbf{10--40$\times$ deeper}
than $\delta$ --- below the diurnal noise, in the quiet zone set by the geothermal
gradient. This underpins \code{MIN\_DEPTH\_CM = 80} and the Anchor Point Method.

\begin{figure}[h]
\centering
\includegraphics[width=0.78\textwidth]{fig_skin_depth.png}
\caption[The daily heat wave dies out with depth]{The size of the day/night
temperature swing as you go down, for both sites (note the \emph{logarithmic}
horizontal axis). \textbf{What to notice:} the swing starts at ${\sim}100$~K at
the surface and collapses by a factor of a thousand within about a metre --- the
dotted line marks the $1/e$ skin depth $\delta$, only a few centimetres. Below
that the curve plunges to essentially nothing. \textbf{Why it matters:} the
Apollo sensors used for the fit all sit well below this skin, in the quiet zone
where temperature is set purely by the deep heat flow.}
\label{fig:skin}
\end{figure}

\subsection{Radiation, sunlight, and the geothermal gradient --- full picture}
\label{sec:boundaries}

\subsubsection{Stefan--Boltzmann: why $T^4$?}
Every surface emits radiation $q_{\mathrm{rad}}=\varepsilon\sigma T^4$. A slightly
hotter surface radiates vastly more ($2\times$ temperature $\Rightarrow 16\times$
power). This makes the surface balance non-linear (Newton solve, \autoref{sec:newton}).

\subsubsection{Surface energy balance}
\begin{equation}
  \underbrace{(1-A)\,S(t)}_{\text{absorbed}}
   \;-\; \underbrace{\varepsilon\sigma T_s^4}_{\text{radiated}}
   \;-\; \underbrace{q_{\mathrm{cond}}(0)}_{\text{into soil}}
   \;=\; 0.
  \label{eq:surface}
\end{equation}
Implemented in \code{surface\_energy\_balance\_residual} (\file{lunar/solver.py}).

\subsubsection{Insolation model}
\[
  S(t) = S_0\cos(\mathrm{lat})\,\max\!\bigl(0,\cos(2\pi t/T_{\mathrm{lunation}})\bigr).
\]
Built in \code{run\_with}. Idealised diurnal cosine; neglected terms are bounded in
the paper.

\subsubsection{Geothermal gradient --- full derivation}
Bottom BC: $-K\,\partial T/\partial z|_{\mathrm{bottom}} = Q_b$.
Time-average the heat equation in periodic steady state ($\partial\langle T\rangle/\partial t=0$):
\[
  \frac{\partial}{\partial z}\!\Bigl[K\,\frac{\partial\langle T\rangle}{\partial z}\Bigr]=0
  \quad\Rightarrow\quad K\,\frac{\mathrm{d}\langle T\rangle}{\mathrm{d}z}=\mathrm{const}=Q_b.
\]
Hence
\begin{equation}
  \frac{\mathrm{d}\langle T\rangle}{\mathrm{d}z} = \frac{Q_b}{K}.
  \label{eq:geotherm}
\end{equation}
Larger $Q_b$ or smaller $K$ $\Rightarrow$ steeper rising mean profile. Example:
$Q_b=0.021$~\si{W/m^2}, $K_d\approx 0.0046$~\si{W/m/K} $\Rightarrow$
gradient $\approx\SI{4.6}{K/m}$.

This \Qb{} matters for the retrieval: deep $T$ depends on the \textbf{ratio}
$Q_b/K_d$ --- the \textbf{\Kd--\Qb{} degeneracy} (\autoref{sec:degeneracy}).

\begin{figure}[h]
\centering
\includegraphics[width=0.78\textwidth]{fig_mean_profile.png}
\caption[Below the skin the mean temperature keeps rising]{The cycle-averaged
temperature $\langle T\rangle$ versus depth for both sites, from the certified
equilibrium solver. \textbf{What to notice:} below ${\sim}1$~m the curve becomes
a straight, \emph{tilted} line --- it does not flatten out. That tilt is the
geothermal gradient $Q_b/K$ of Eq.~\eqref{eq:geotherm}. \textbf{Why it matters:}
the deep slope encodes the ratio $Q_b/K_d$; reading the Apollo temperatures
against this line is exactly how the project pins down \Kd.}
\label{fig:meanprofile}
\end{figure}

\subsection{Boundary conditions --- full derivation and connection to the main equation}
\label{sec:bc-full}

\subsubsection{Why the main equation is not enough}
The heat equation (Eq.~\eqref{eq:heat}) is second-order in $z$, so it needs
\textbf{two boundary conditions} plus an initial condition $T(z,t{=}0)$ to pick a
unique solution. The complete problem: interior PDE $+$ top BC $+$ bottom BC $+$
initial profile (paper Eqs.~1--3).

\subsubsection{One unifying idea: energy balance at every control volume}
Interior cells: flux in minus flux out equals storage. At boundaries the same
bookkeeping applies, but one face is the domain edge (space above, or deep
interior below) --- not another soil cell.

\subsubsection{Top boundary --- radiative surface balance}
At $z=0$ there is no conduction from above. Energy balance on the surface skin:
\begin{equation}
  (1-A)\,S(t) \;-\; \varepsilon\sigma T_s^4 \;-\; q_{\mathrm{cond}}(0) \;=\; 0.
  \label{eq:bc-top-teach}
\end{equation}
The conductive term is Fourier's law at the top face:
$q_{\mathrm{cond}}(0)=-K\,\partial T/\partial z|_{0+}\approx K_{\mathrm{surf}}(T_s-T_{\mathrm{sub}})/(\Delta z_0/2)$,
implemented in \code{surface\_energy\_balance\_residual}. Each timestep,
\code{\_solve\_surface\_newton} finds $T_s$; the first row of the tridiagonal
system uses $T_s$ as a ghost temperature coupling the BC to the interior PDE.

\subsubsection{Bottom boundary --- geothermal flux}
Below the 5~m column we prescribe upward flux $Q_b$ (Neumann condition):
\begin{equation}
  -\,K\,\frac{\partial T}{\partial z}\Big|_{z=z_{\mathrm{max}}} \;=\; Q_b.
  \label{eq:bc-bot-teach}
\end{equation}
For the deepest cell, the outgoing flux at the bottom face is fixed to $Q_b$
rather than computed from a neighbour. In \code{\_step}:
\code{d[-1] += dt * Q\_b / cap[-1]} adds the basal energy source. In steady
state this forces $K\,\mathrm{d}\langle T\rangle/\mathrm{d}z=Q_b$ everywhere
(Eq.~\eqref{eq:geotherm}).

\subsubsection{How the pieces fit in one timestep}
\begin{enumerate}[leftmargin=1.5em,itemsep=2pt]
  \item \textbf{Top:} Newton solve for $T_s$ from radiative balance.
  \item \textbf{Interior:} Crank--Nicolson discretisation of Eq.~\eqref{eq:heat}.
  \item \textbf{Bottom:} add $Q_b$ source to the last cell.
  \item \textbf{Coupled solve:} one Thomas tridiagonal sweep for $T^{n+1}$.
\end{enumerate}
The BCs are not separate from the PDE --- they are the same energy-conservation
principle at the domain edges, wired into the same linear system.

\begin{headsup}
\textbf{Sanity check.} Day/night surface swings $\leftarrow$ top BC;
rising mean profile $\leftarrow$ bottom BC; quiet deep sensors $\leftarrow$
interior damping (skin depth). Wrong BCs would still run but would not match
Apollo or Diviner.
\end{headsup}

\subsection{Orbital mechanics and the top boundary --- full picture with diagrams}
\label{sec:orbit-top-bc}

The \textbf{top boundary} is where the Moon meets space. Unlike Earth, there is
no atmosphere to scatter or absorb sunlight, and no wind to carry heat away.
Everything that heats or cools the surface comes from:
\begin{enumerate}[leftmargin=1.4em,itemsep=2pt]
  \item \textbf{Direct sunlight} (orbital geometry sets how strong it is at each moment),
  \item \textbf{Thermal radiation} the surface emits to the cold sky,
  \item \textbf{Conduction} into the regolith below.
\end{enumerate}
This subsection explains the \textbf{orbital part} (where does $F_\odot(t)$ come
from?) and ties it to the \textbf{surface energy balance} that closes the top BC.

\subsubsection{The Moon in orbit: what ``one day'' means}

The Moon rotates once per orbit around Earth, so the same face always points
toward Earth. From a point on the lunar surface, ``day'' and ``night'' are set by
the Sun, not by watching Earth.

\begin{itemize}[leftmargin=1.4em,itemsep=2pt]
  \item \textbf{Synodic month} $P \approx 29.53$~days: time from one local noon
    (Sun highest) to the next. This is the period used in the code
    (\code{T\_LUNAR} $\approx 2.55\times10^6$~\si{s}).
  \item \textbf{Subsolar point}: the spot on the Moon where the Sun is directly
    overhead ($\theta_\odot = 0$). It marches around the equator once per synodic month.
  \item \textbf{Site latitude} $\phi$: Apollo~15 is at $26.1^\circ$N, Apollo~17 at
    $20.2^\circ$N. The Sun never passes directly overhead there; at local ``noon'' the
    Sun is lower in the sky than at the equator.
\end{itemize}

\begin{plainenglish}
Imagine standing on the Moon at an Apollo site. The Sun rises, tracks across the
sky at a height that depends on your latitude, sets, and you wait through a long
night (${\sim}$14 Earth days) until it rises again. The model does not simulate
every crater shadow or mountain --- it uses a smooth ``average day'' for that latitude.
\end{plainenglish}

% ---- Figure: Sun--Moon--surface geometry --------------------------------
\begin{figure}[htbp]
\centering
\begin{tikzpicture}[
  x=1cm, y=1cm,
  font=\small,
  sun/.style={circle, fill=yellow!80!orange, draw=orange!70!black, line width=0.6pt, minimum size=1.1cm},
  moon/.style={circle, fill=softgrey!40, draw=dim, line width=0.8pt},
  ray/.style={-{Stealth[length=2.5mm]}, line width=0.7pt, color=orange!85!black},
  dashedline/.style={dashed, color=dim, line width=0.5pt},
  label/.style={font=\footnotesize, text=char}
]

% Sun
\node[sun] (sun) at (-5.5, 2.2) {};
\node[label, below=2pt of sun] {Sun};

% Moon (cross-section)
\draw[moon] (0,0) circle (2.2);
\fill[softgrey!25] (0,0) circle (2.2);

% Local horizontal at site (latitude phi)
\coordinate (C) at (0,0);
\coordinate (S) at ({2.2*cos(55)},{2.2*sin(55)});  % site on surface, phi ~ 35 deg from equator visually
\draw[line width=1.2pt, color=char] (C) -- (S);
\node[label, anchor=west] at ($(S)+(0.15,0.05)$) {regolith surface};

% Vertical at site (local zenith)
\draw[dashedline] (S) -- ++(0, 1.8) node[above, label] {local vertical (zenith)};

% Sun ray to site
\draw[ray] (sun.east) .. controls (-2.5, 2.0) and (-0.8, 1.6) .. (S);
\node[label, fill=paper, inner sep=1pt] at (-2.2, 2.55) {sunlight};

% Solar zenith angle theta
\draw[{Stealth[length=2mm]}-, line width=0.6pt, color=teal] (0.35, 0) arc (0:55:0.55);
\node[label, text=teal] at (0.95, 0.22) {$\theta_\odot$};

% Latitude phi from equator
\draw[dashedline] (-2.2, 0) -- (2.2, 0);
\draw[{Stealth[length=2mm]}-, line width=0.6pt, color=forest] (0.7, 0) arc (0:55:0.7);
\node[label, text=forest] at (1.35, -0.15) {$\phi$};

% Equator hint
\node[label, text=dim] at (-2.5, -0.35) {lunar equator (plane of subsolar march)};

% Orbit hint
\draw[-{Stealth[length=2mm]}, line width=0.5pt, color=plum, dashed]
  (2.5, 0.3) arc (10:350:2.5 and 0.35);
\node[label, text=plum, anchor=west] at (2.9, 0.3) {Moon orbit (schematic)};

% Annotation box
\node[draw=teal!60, fill=teal!5, rounded corners=2pt, align=left, font=\footnotesize,
      text width=4.2cm, anchor=north west] at (-5.8, -2.6) {
  \textbf{Zenith angle} $\theta_\odot$: angle between\\
  local vertical and the Sun.\\
  Smaller $\theta_\odot$ $\Rightarrow$ stronger heating.
};

\end{tikzpicture}
\caption{\textbf{Sun--Moon--surface geometry.} At a site at latitude $\phi$, the
  Sun is not overhead at ``noon'' unless $\phi=0$. The \textbf{solar zenith angle}
  $\theta_\odot$ sets how much solar power hits each square metre of horizontal
  ground (cosine law). The retrieval uses an idealised cycle
  $\cos\theta_\odot(t)=\cos\phi\cos(2\pi t/P)$; full SPICE ephemeris lives in
  \file{lunar/ephem.py} but is not required for the \Kdstar{} fit.}
\label{fig:sun-moon-geometry}
\end{figure}

\subsubsection{From geometry to insolation: the cosine law}

The solar constant $S_0 \approx \SI{1361}{W/m^2}$ is the flux in a beam
perpendicular to the Sun. On a \textbf{horizontal} patch of ground, only the
component normal to the beam counts:
\begin{equation}
  F_\odot(t) = \max\!\bigl(0,\; S_0 \cos\theta_\odot(t)\bigr).
  \label{eq:insolation-teach}
\end{equation}
When the Sun is low ($\theta_\odot$ large), $\cos\theta_\odot$ is small and
heating is weak. At night, $\cos\theta_\odot < 0$ and the $\max(0,\cdot)$ clips
insolation to zero.

\paragraph{Idealised synodic model used in production.}
For each site latitude $\phi$,
\begin{equation}
  \cos\theta_\odot(t) = \cos\phi \,\cos\!\left(\frac{2\pi t}{P}\right),
  \label{eq:cos-zenith-ideal}
\end{equation}
with $P = \code{T\_LUNAR}$. This is the same formula as in the manuscript
(Eq.~\eqref{eq:insolation-teach}). It assumes:
\begin{itemize}[leftmargin=1.4em,itemsep=1pt]
  \item smooth variation over one lunation (no libration jitter at the hour scale),
  \item no solar \textbf{declination} tilt ($\pm 1.5^\circ$ neglected),
  \item no orbital \textbf{eccentricity} changing $S_0$ ($<0.1\%$ effect on cycle mean).
\end{itemize}
Those corrections matter for \emph{phase-resolved} work; this project compares
\textbf{time-mean} temperatures over a stability window, so the simplified cycle is
adequate (see manuscript \S``Surface solar geometry'').


% ---- Figure: insolation over one lunation --------------------------------
\begin{figure}[htbp]
\centering
\begin{tikzpicture}[x=0.42cm, y=0.00035cm, font=\small]
  % Axes
  \draw[-{Stealth}, line width=0.6pt] (0,0) -- (52,0) node[right, font=\footnotesize] {$t$ (days)};
  \draw[-{Stealth}, line width=0.6pt] (0,0) -- (0,4200)
    node[above, font=\footnotesize] {$F_\odot$ (\si{W/m^2})};
  \draw[color=dim!60, line width=0.3pt] (0,0) -- (0,4200);
  \draw[color=dim!60, line width=0.3pt] (7.38,0) -- (7.38,4200);
  \draw[color=dim!60, line width=0.3pt] (14.76,0) -- (14.76,4200);
  \draw[color=dim!60, line width=0.3pt] (22.14,0) -- (22.14,4200);
  \draw[color=dim!60, line width=0.3pt] (29.53,0) -- (29.53,4200);
  \node[below, font=\tiny, text=dim] at (0,0) {0};
  \node[below, font=\tiny, text=dim] at (7.38,0) {7.4};
  \node[below, font=\tiny, text=dim] at (14.76,0) {14.8};
  \node[below, font=\tiny, text=dim] at (22.14,0) {22.1};
  \node[below, font=\tiny, text=dim] at (29.53,0) {29.5};
  \node[below, font=\footnotesize, text=dim] at (26, -28)
    {one synodic month $P \approx 29.5$~d};

  % cos curve for phi = 26 deg (peak ~ 1223 W/m2); day half only
  \draw[line width=1.1pt, color=coral, domain=0:14.76, samples=80]
    plot (\x, {1361*0.899*cos(360*\x/29.53)});

  % Night shading
  \fill[dim!12] (14.76,0) rectangle (29.53,4200);

  \node[font=\footnotesize, text=dim, rotate=90] at (-3.5, 600) {night ($F_\odot=0$)};
  \node[font=\footnotesize, text=coral] at (7.5, 1150) {daytime insolation};
  \node[font=\footnotesize, align=center] at (38, 3500) {Apollo~15-scale\\$\phi=26.1^\circ$N\\peak $\lesssim S_0\cos\phi$};

  % S0 reference
  \draw[dashed, color=teal, line width=0.6pt] (0,1361) -- (50,1361);
  \node[font=\tiny, text=teal, anchor=west] at (50.5, 1361) {$S_0$};

\end{tikzpicture}
\caption{\textbf{Idealised insolation} $F_\odot(t)$ for northern mid-latitude
  ($\cos\phi \approx 0.90$). Sunlit half of the lunation: flux follows a cosine;
  night: clipped to zero. Peak is below $S_0$ because the Sun never reaches the
  zenith at $\phi > 0$. Code: \code{insol = S0 * cos\_lat * np.maximum(0, np.cos(phase))}
  in \file{retrieve\_kd.py}.}
\label{fig:insolation-lunation}
\end{figure}

\subsubsection{Top boundary: surface energy balance (control volume)}

Take a thin control volume \emph{at} the surface --- the top face sees space, the
bottom face touches the first regolith cell. In steady balance at each instant,
Eq.~\eqref{eq:bc-top-teach} applies (derived in \autoref{sec:bc-full}), with
insolation written explicitly as $F_\odot(t)$ instead of $S(t)$:
\[
  (1-A)\,F_\odot(t) \;=\; \varepsilon\sigma T_s(t)^4 \;+\; q_{\mathrm{cond}}(t).
\]
Here $q_{\mathrm{cond}} = -K\,\partial T/\partial z\big|_{z=0}$ is the downward
conductive flux (positive when heat enters the ground).

% ---- Figure: surface energy balance --------------------------------------
\begin{figure}[htbp]
\centering
\begin{tikzpicture}[
  x=1cm, y=1cm, font=\small,
  arr/.style={-{Stealth[length=2.8mm]}, line width=0.85pt},
  thinarr/.style={-{Stealth[length=2mm]}, line width=0.6pt}
]

% Regolith column
\fill[softgrey!35] (-1.2, -2.8) rectangle (1.2, 0);
\draw[line width=0.8pt, color=char] (-1.2, -2.8) rectangle (1.2, 0);
\foreach \y in {-0.5,-1.0,-1.5,-2.0,-2.5} {
  \draw[dim!40, line width=0.3pt] (-1.2,\y) -- (1.2,\y);
}
\node[font=\footnotesize, text=dim] at (0, -1.4) {regolith ($z>0$)};

% Surface line
\draw[line width=1.4pt, color=char] (-1.5, 0) -- (1.5, 0);
\node[anchor=west, font=\footnotesize] at (1.55, 0) {$z=0$ (surface)};

% Sun / insolation from upper left
\draw[arr, color=orange!85!black] (-2.8, 2.2) -- (-0.3, 0.35);
\node[align=left, font=\footnotesize, fill=paper, inner sep=2pt] at (-3.5, 2.5) {
  incident solar\\$F_\odot(t)$
};
\draw[thinarr, color=orange!60!black] (-1.0, 1.5) -- (-0.2, 0.25);
\node[align=left, font=\footnotesize, text=orange!70!black] at (-2.2, 1.0) {
  absorbed\\$(1-A)F_\odot$
};

% Reflected (small arrow back)
\draw[thinarr, color=dim, dashed] (-0.1, 0.2) -- (-1.2, 1.0);
\node[font=\tiny, text=dim] at (-1.5, 1.15) {reflected $AF_\odot$};

% Outgoing IR upward
\draw[arr, color=coral] (0.2, 0.15) -- (0.2, 2.0);
\node[align=left, font=\footnotesize, text=coral, anchor=west] at (0.45, 2.0) {
  emitted thermal radiation\\$\varepsilon\sigma T_s^4$
};

% Conduction downward
\draw[arr, color=teal] (0, -0.05) -- (0, -1.2);
\node[align=center, font=\footnotesize, text=teal, anchor=west] at (0.25, -0.7) {
  conduction\\$q_{\mathrm{cond}}=-K\partial T/\partial z$
};

% Sky
\fill[teal!4] (-3.2, 0.05) rectangle (3.2, 2.8);
\node[font=\footnotesize, text=dim] at (0, 2.45) {space / sky ($T_{\mathrm{space}}\approx 0$)};

% Ts label
\node[circle, fill=coral!15, draw=coral, inner sep=1pt, font=\footnotesize\bfseries] at (0, 0) {$T_s$};

% Balance equation box
\node[draw=char!50, fill=paper, rounded corners=3pt, align=center, font=\footnotesize,
      text width=5.8cm] at (0, -3.55) {
  Top BC (each timestep): \quad
  $(1-A)F_\odot = \varepsilon\sigma T_s^4 + q_{\mathrm{cond}}$\\[3pt]
  \code{surface\_energy\_balance\_residual(Ts)} = 0
  \quad$\Rightarrow$\quad Newton finds $T_s$
};

\end{tikzpicture}
\caption{\textbf{Top boundary control volume.} Absorbed sunlight (after albedo $A$)
  must equal the sum of infrared emission and heat conducted into the regolith.
  $T_s$ is not stored as a fixed number --- it is \textbf{solved} each hour so the
  balance holds. This is what \code{\_solve\_surface\_newton} does in
  \file{solver.py}.}
\label{fig:surface-energy-balance}
\end{figure}

\begin{analogy}
The surface is like a pan on a stove that also glows in the dark. Sunlight pours
in (when it's day), the pan radiates heat to the cold kitchen (space), and the
handle (regolith) conducts heat downward. At every instant the pan's top temperature
adjusts until those three flows balance.
\end{analogy}

\subsubsection{How orbital forcing connects to the heat equation}

The chain from orbit to interior is:
\begin{center}
\begin{tikzpicture}[
  node distance=0.35cm and 0.5cm,
  box/.style={draw, rounded corners=2pt, fill=paper, align=center,
              font=\footnotesize, inner sep=5pt, minimum width=2.6cm},
  arr/.style={-{Stealth[length=2mm]}, line width=0.55pt, color=char}
]
\node[box, fill=orange!8] (orbit) {Orbit + latitude $\phi$\\$\Rightarrow\cos\theta_\odot(t)$};
\node[box, fill=coral!8, right=of orbit] (insol) {$F_\odot(t)=S_0\max(0,\cos\theta_\odot)$};
\node[box, fill=teal!8, right=of insol] (bc) {Top BC:\\$(1-A)F_\odot=\varepsilon\sigma T_s^4+q_{\mathrm{cond}}$};
\node[box, fill=forest!8, right=of bc] (pde) {Interior PDE:\\$\rho c_p\partial T/\partial t=\partial(K\partial T/\partial z)/\partial z$};
\node[box, fill=plum!8, below=of bc] (newton) {Newton $\to T_s$ each step};
\draw[arr] (orbit) -- (insol);
\draw[arr] (insol) -- (bc);
\draw[arr] (bc) -- (pde);
\draw[arr] (bc) -- (newton);
\draw[arr] (newton) |- (pde);
\end{tikzpicture}
\end{center}

Insolation is \textbf{prescribed} (known input from orbit model). Skin temperature
$T_s$ is \textbf{unknown} until the top BC is satisfied. Once $T_s$ is known, it
feeds the first interior cell and Crank--Nicolson advances $T(z,t)$ for all deeper
layers.

\subsubsection{Bottom boundary (for completeness)}

\begin{figure}[htbp]
\centering
\begin{tikzpicture}[x=1cm, y=1cm, font=\small,
  arr/.style={-{Stealth[length=2.5mm]}, line width=0.75pt}]
\fill[softgrey!35] (-1.2, 0) rectangle (1.2, 2.5);
\draw[line width=0.8pt] (-1.2, 0) rectangle (1.2, 2.5);
\node[font=\footnotesize, text=dim] at (0, 1.2) {regolith};
\draw[line width=1.2pt] (-1.5, 0) -- (1.5, 0);
\node[anchor=west, font=\footnotesize] at (1.55, 0) {$z=z_{\mathrm{max}}$ (5~m)};
\draw[arr, color=coral] (0, -0.1) -- (0, -1.1);
\node[align=center, font=\footnotesize, text=coral, anchor=west] at (0.2, -0.65) {
  basal flux $Q_b$\\from lunar interior
};
\node[draw=char!40, fill=paper, rounded corners=2pt, font=\footnotesize, align=center,
      text width=5.5cm] at (0, -2.0) {
  Bottom BC: \quad $-K\dfrac{\partial T}{\partial z}\Big|_{z_{\mathrm{max}}} = Q_b$\\[2pt]
  Weak interior heat; ${\sim}0.02$~\si{W/m^2} at Apollo sites
};
\end{tikzpicture}
\caption{\textbf{Bottom boundary.} A small steady flux from the Moon's interior.
  No orbital variation --- only the top sees the Sun.}
\label{fig:bottom-bc}
\end{figure}

\begin{headsup}
\textbf{What is \emph{not} in the top BC.} No atmosphere, no wind, no latent heat
(evaporation), no ground heat storage in a separate surface layer. The model is
deliberately minimal: radiative balance + 1-D conduction. That matches airless-body
practice for metre-scale regolith (Hayne et al.\ 2017; Vasavada et al.\ 2012).
\end{headsup}

\subsubsection{Code map: from \code{insolation} to \code{\_step}}

\begin{table}[htbp]
\centering
\small
\begin{tabular}{@{}llp{6.2cm}@{}}
\toprule
Step & Function / variable & Role \\
\midrule
1 & \code{cos\_lat = np.cos(np.radians(lat))} & latitude factor \\
2 & \code{phase = 2*np.pi*t/T\_LUNAR} & lunation phase \\
3 & \code{insol = S0 * cos\_lat * max(0, cos(phase))} & build $F_\odot(t)$ \\
4 & \code{run\_with(..., insolation=insol)} & pass into solver \\
5 & \code{\_solve\_surface\_newton(T\_prev, insol\_k, ...)} & top BC $\to T_s$ \\
6 & \code{\_step(...)} & CN step for interior + bottom $Q_b$ \\
\bottomrule
\end{tabular}
\caption{End-to-end path for orbital forcing and boundary conditions in the retrieval pipeline.}
\label{tab:orbit-bc-code}
\end{table}

For full SPICE-based zenith angles (sub-solar latitude, libration, exact timestamps),
see \file{lunar/ephem.py}. The \Kdstar{} retrieval does not require that path because
it fits \emph{mean} temperatures, not diurnal phase.

